\section{Conclusions}
\label{sec:conclusions}

This study used DEM simulations of undrained cyclic torsional shear in hollow cylindrical specimens to investigate how initial stress anisotropy (\Kzero{}) affects liquefaction resistance under HCA-type boundary conditions. The main findings are summarized as follows:

\begin{enumerate}
\item A coupled servo formulation was developed for the DEM-HCA model. The pressure differences ($p_{dif,r}'$ and $p_{dif,z}'$) were identified as the correct servo targets for mapping laboratory boundary conditions, and the constant-volume constraint was incorporated to couple the radial and axial degrees of freedom. Analytical inversion of the resulting stiffness matrix yields closed-form servo coefficients that ensure numerical stability.

\item Liquefaction resistance exhibits a monotonic directional dependency on initial stress anisotropy: compression-state anisotropy (\Kzero{} $< 1.0$) produces higher resistance, while extension-state anisotropy (\Kzero{} $> 1.0$) reduces it. This ordering is consistent across all CSR levels and both relative densities examined. The per-cycle energy input rate differs because of distinct stiffness degradation trajectories: the compression-state specimen degrades stiffness most slowly, and the extension-state specimen fastest.

\item The mechanical coordination number \Zm{} and the signed fabric anisotropy indicator $\alpha$ play complementary roles. \Zm{} primarily reflects packing compactness and separates the two density groups, but does not independently explain the \Kzero{}-dependent resistance within each group. By contrast, $\alpha_0$ correlates positively with $N_L$ within both density groups, indicating that directional fabric governs the resistance variation among \Kzero{} states. Decomposition of the fabric tensor in cylindrical coordinates shows that a contact's contribution to torsional resistance depends on the projection of its normal onto the $z$--$\theta$ plane; compression-state consolidation concentrates a larger fraction of contacts onto this plane, offering a geometric interpretation of the $\alpha_0$--$N_L$ correlation. Both the count-weighted and force-weighted fabric tensors capture this directional signal through their torsion-resisting fractions ($1 - \Phi_{rr}$ and $1 - \Phi^f_{rr}$).

\item Multiple indicators, including shear wave velocity, mechanical coordination number, signed fabric anisotropy indicator, count- and force-weighted fabric tensors, contact-density distributions, and force-chain configurations, all show that initially distinct fabric states progressively degrade and converge toward similar configurations in the post-liquefaction stage. This convergence indicates that cyclic loading drives the granular assembly into a common dynamic equilibrium in which the fabric oscillates with loading reversals, regardless of the initial stress anisotropy.

\end{enumerate}
